%! Polynomial Functions and Rates of Change — Unit 1, Lesson 1.4 — Exit Ticket (Answer Key)
\documentclass[10pt]{article}
\usepackage{precalculus-article}
\usepackage{precalculus-key}   % pulls in -boxes; do NOT also load -boxes
\newcommand{\TallMath}[1]{$\displaystyle #1\rule[-1.4em]{0pt}{3.2em}$}

\begin{document}

\pageheader{Unit 1, Lesson 1.4}{Exit Ticket --- Answer Key}
\namedateperiod

\vspace{0.15in}

\begin{notesbox}{Exit Ticket --- Key}
\begin{enumerate}[leftmargin=1.5em, itemsep=10pt]

\item A polynomial function $p$ has exactly three distinct real zeros.

\begin{enumerate}[leftmargin=1.5em, itemsep=6pt]
    \item[(a)] Minimum number of local extrema:

    Answer: \ans{2}

    \item[(b)] Supporting principle:

    \ans{By EK 1.4.A.3, between every two distinct real zeros of a nonconstant polynomial there must be at least one local maximum or minimum. With three zeros there are two consecutive pairs, each requiring at least one local extremum --- so at least 2 local extrema total.}
\end{enumerate}

\item A polynomial function $q$ has a point of inflection at $x = 2$.
The AROC of $q$ over $[0, 2]$ is 3, and the AROC of $q$ over $[2, 4]$ is 7.

\begin{enumerate}[leftmargin=1.5em, itemsep=6pt]
    \item[(a)] Is the AROC increasing or decreasing?

    \ans{Increasing} --- the AROC goes from 3 to 7.

    \item[(b)] Concavity and explanation:

    \ans{Concave up. When the AROC is increasing as $x$ increases, the graph is concave up. The rate of change switches from a smaller value (3) to a larger value (7) as we pass through $x=2$, which is consistent with the graph curving upward. The point of inflection at $x=2$ is where this switch occurs.}
\end{enumerate}

\end{enumerate}
\end{notesbox}

\end{document}
